% Source-status section, audited 12 October 2026.
% Citation keys are supplied in literature-references.bib.
\section{Sources and the announced quasipolynomial bound}
\label{sec:literature-status}

The source record contains two different quantitative versions of Lackenby's
2021 announcement.  They should not be silently identified.
\begin{center}
\begin{tabular}{lll}
\hline
Source & Stated running time & Stated hierarchy depth\\
\hline
February 2021 author slides & $n^{O(\log n)}$ & $O(\log n)$\\
March 2021 Oxford slides & $2^{O((\log n)^3)}$ & $O((\log n)^2)$\\
15 April 2021 Princeton abstract & $2^{O((\log n)^3)}$ & not specified\\
\hline
\end{tabular}
\end{center}
Here $n$ is the number of crossings of the input knot diagram.  The first
row is explicit in the February slides, including their complexity summary;
it also agrees with Oxford's news announcement
\cite{Lackenby2021February,Oxford2021Announcement}.  The later slide version
and Princeton abstract state the weaker, still quasipolynomial bound
\cite{Lackenby2021Oxford,Lackenby2021Princeton}.  Thus the earlier
$n^{O(\log n)}$ figure is a documented historical claim, while
\[
  2^{O((\log n)^3)}=n^{O((\log n)^2)}
\]
is the later explicit benchmark used here.  Matching the stronger exponent
would require an additional argument, not a change of notation.

The later slides describe a hierarchy procedure whose cost is controlled by
its depth and a polynomially bounded surface parameter.  They combine
compressed parallelity bundles, multisurfaces, and modifications of
generalised Heegaard splittings.  Their logarithmic-squared depth bound uses
a dichotomy involving boundary-genus growth and special regions satisfying
a Cheeger-type inequality \cite{Lackenby2021Oxford}.  This architecture is
more than an efficient enumeration of fundamental normal surfaces.  A
multisurface in those slides consists of individually embedded surfaces
representing independent relative homology classes; the surfaces may
intersect, and successive cuts are taken in their complements.  A collection
of simultaneously disjoint representatives cannot be substituted without
proof.

As of this article's source audit on 12 October 2026, the sources located for
the quasipolynomial complexity theorem are announcements and slides.  A
complete published proof of that bound was not located.  This is a statement
about the material found in the audit, not a claim that no other manuscript
exists.  The author posted a substantial new preprint,
\emph{Incompressible surfaces, hierarchies and unknot recognition}, on
25 July 2026 \cite{Lackenby2026Hierarchies}.  It develops an algorithm for
incompressibility and essential boundary patterns in irreducible compact
orientable $3$-manifolds, and discusses future speed-ups.  Proposition~9.1
bounds an iteration count by $L(g+1)^L$ in terms of the maximum hierarchy
length $L$ and pattern complexity $g$; it does not supply the bounds on those
parameters needed for the announced exponent.  Theorem~14.4 provides
compressed parallelity-bundle extraction for uniform handle structures.
Its stronger admissibility procedure in Proposition~14.2 assumes both
irreducibility and an essential boundary pattern.  Those hypotheses must be
checked before using that procedure in an unknot-recognition implementation.

Several necessary ingredients are already published.  Agol, Hass, and
Thurston give polynomial-time orbit counting for interval pairings and a
weighted extension useful for normal-surface component data
\cite{AgolHassThurston2006}.  Their weighted output permits an interval of
orbits with one common weight vector; multiplicities need not be expanded.
This matters whenever the normal input contains exponentially many repeated
components.  Lackenby's certification paper supplies the explicit core and
ordered interval-bundle attachments used in
Section~\ref{sec:unmarked-reconstruction}
\cite{Lackenby2021Certification}.  These results justify compressed cutting,
but do not establish a fast method of finding a suitable short hierarchy.

Boundary markings have their own algorithmic theory.  The published
\emph{Some Fast Algorithms for Curves in Surfaces} gives algorithms for
geometric intersection numbers and, for essential curves and arcs, isotopy,
with polynomial dependence on triangulation size and logarithmic dependence
on normal weights \cite{Lackenby2026Curves}.  This offers a concrete source of
subroutines for a future marking-aware implementation.  It does not by itself
bound the complexity of transferring an arbitrary hierarchy pattern through
a $3$-dimensional reconstruction.  The research program must state that
transfer problem separately and prove its complexity.
